\documentclass[11pt]{article}
\usepackage[margin=1in]{geometry}
\usepackage{amsmath,amssymb,hyperref}
\newcommand{\F}{\mathbb F}
\newcommand{\E}{\mathbb E}
\newcommand{\Res}{\operatorname{Res}}
\newcommand{\wtO}{\widetilde O}
\title{Implicit halves for the scalar resultant}
\author{Project research note}
\date{18 November 2026}
\begin{document}
\maketitle

The normalized norm and its first derivatives admit an exact evaluation
in which neither of the two original half polynomials is constructed.
Three smaller blocks suffice. At the tuned widths this gives a
time--space tradeoff on the relation queries. The space exponent
drops from $1/2$ to $1/3$ in units of $b$, and the time exponent
increases from $1/2$ to $2/3$. A separate algorithm computes a
small-degree gcd probe on a genuine hashed slot in sublinear work
in its width. We examine what such probes and rational-series
compression can reveal about the final scalar.

\section{A symmetric three-block evaluation}
Use the smooth curve $y^2=x^3+ax+c$, $p>3$, the fixed target
$2T\ne O$, and the normalized sign-orbit form $S$ of
\texttt{../scalar\_norm/norm.tex}. With $t=x(T)$ and paired
abscissa slots $U_i$, the requested scalar and derivatives are
those of
\[
 N(z,\mathbf s)=
 \prod_{\mathbf u\in\prod U_i}\left(S(\mathbf u;z)
                     +s_0+\sum_i s_i u_i\right)
\]
at $(t,\mathbf0)$. The sign-orbit construction uses the triquadratic
\[
 K(x,y,z)=(x-y)^2z^2-2((x+y)(xy+a)+2c)z+(xy-a)^2-4c(x+y),
\]
homogenized in each argument. It accounts for both completions
$x(P+Q),x(P-Q)$ including infinity and multiplicities.

Include $\{t\}$ as one additional singleton slot. Partition these
$m+1$ sets into three nonempty groups. For each group of $s$ sets,
multiply the individual forms $S(\mathbf u;X)$ for its tuples.
The resulting form has formal degree $2^{s-1}\prod n_i$, where
$n_i$ is the size of each set and the target has size one. Denote
the three forms by $F,G,H$ with degrees $d_F,d_G,d_H$. In particular
\begin{equation}
    d_F d_G d_H=2^{m-2}\prod_i|U_i|=\tfrac14\prod_i|F_i|,
    \label{eq:volume}
\end{equation}
where $F_i$ denotes the signed point slot. One may interleave the
three groups across the original two halves. The new evaluator
requires neither original half nor a list of its tuple roots.

Factor a stationary block $F$ into normalized finite linear factors
$X_0-xX_1$ and infinity factors $X_1$. Let $\alpha_F$ be their
common scalar prefactor, the highest nonzero affine coefficient.
List its projective roots as $\xi$ with multiplicities $m_\xi$.
Use the same conventions for $G$. Holding these two blocks fixed, form
one polynomial per distinct outer root,
\[
 P_\xi(Z)=\prod_{\eta:G(\eta)=0}K(\xi,\eta,Z)^{m_\eta},
                  \qquad \deg_{\rm formal}P_\xi=2d_G ,
\]
where the product indexes distinct projective roots. The exact value is
\begin{equation}
 N(t,\mathbf0)=
 \alpha_F^{\,2d_Gd_H}\alpha_G^{\,2d_Fd_H}
       \prod_{\xi:F(\xi)=0}
       \Res_{2d_G,d_H}(P_\xi,H)^{m_\xi}.             \label{eq:three}
\end{equation}
This is the Poisson product formula for three successive homogeneous
eliminations against $K$. To verify the normalization, first use
one tuple per block. Merging two forms gives one $K$ factor for
every pair of sign orbits. Its roots represent the two possible
relative signs. This is the homogeneous resultant recurrence in
the preceding note, reparenthesized by multiplicativity. The
remaining elimination pairs these relative signs with the third
group. The full normalized summation form is symmetric in all
$m+1$ arguments for $m\ge2$, as proved there. Taking products of
the tuple identities gives (\ref{eq:three}). Each input to $K$
has degree two, so the Poisson signs from infinity factors are
positive, as is the sign for interchanging $H$ and $P_\xi$.
In particular $K(\infty,\eta,Z)=(Z_0-\eta Z_1)^2$ for finite
$\eta$, and $K(\infty,\infty,Z)=Z_1^2$.

To differentiate an input belonging to $H$, differentiate only
$H$ inside (\ref{eq:three}). The two prefactors are fixed.
Each section uses the fixed-form resultant derivative in the
preceding note, which handles a common root and a drop in affine
degree. Apply the ordinary product and power rules, including at
zero. Permute the three groups to cover all inputs. This avoids
differentiating a factorization at a multiple root. The case of
one original slot uses its defining product.

In detail, this gives $N_z$ and the derivatives $P_i,Q_i$ of the
norm under $u\mapsto u+\epsilon$ and $u\mapsto(1+\epsilon)u$ in
slot $i$. At an isolated tuple the previous result gives
$u_i^*=Q_i/P_i$, $N_{s_0}=N_z/S_z(\mathbf u^*;t)$ and
$N_{s_i}=u_i^*N_{s_0}$. At a multiple zero, use three-block
scalar calls on restricted boxes to locate one zero tuple and
obtain its cofactor from the disjoint complement boxes. With
several zero tuples that cofactor vanishes. At a nonzero norm the
relation query can exit immediately. If all additive derivatives
are requested there, the existing two-slot symbolic calculation
computes them without either original half. These two types of
derivatives have different costs. Empty slots and rational
two-torsion use the earlier unit and fixed-point branches.

\paragraph{Cost.}
Build block polynomials by subproduct trees. The roots of each
individual factor can be enumerated by sign sums on the curve or
factored in its small degree $2^{O(m)}$. Build each $P_\xi$ from
quadratic factors with fast polynomial products. With balanced
groups, tuned independent widths and growing polylogarithmic $m$,
all three degrees are at most $2^{b/3+O(b/m+m)}$. One varying-block
pass takes $\wtO(d_F(d_G+d_H))$ field operations times
$2^{O(m)}\operatorname{poly}(b,m)$ allowing for construction,
directions and root enumeration. Peak working space is at most
$2^{O(m)}(d_F+d_G+d_H)\operatorname{poly}(b,m)$. Neither bound
uses an oracle for an uncharged polynomial evaluation.
Thus on the tuned odd prime-order family the relation query, including the first derivatives at zeros,
costs $2^{(2/3+o(1))b}$ time and $2^{(1/3+o(1))b}$ space.
The complete additive gradient at nonzeros has the separate
$2^{(1-1/m)b+O(m+\log b)}$ upper bound of the two-slot method.
Its working polynomial degree is $O(2^m\max_i n_i)$.

Equation (\ref{eq:volume}) prices the extra scan. In the method
that scans all formal outer roots and processes the two remaining
explicit polynomials for each section, work times space is at
least of order $d_Fd_Gd_H$. With fixed balanced target-free blocks
the previous collision estimate also shows that a $1-o(1)$ fraction
of their formal roots are distinct in probability. During the pass
varying the target block, merging repeated outer or middle roots
therefore does not change the exponent. These estimates concern
this evaluation schedule. They do not price a simultaneous
implicit evaluation of all sections by a different identity.

\section{An actual hash probe of sublinear cost}
For modulus $M=2^b$ the canonical candidate abscissas are
$x_j=(A j+C)\bmod M<p$, $0\le j<\tau$, where $A$ is the inverse
of the odd hash multiplier. Write $g_{\rm raw}(X)=\prod_j(X-x_j)$
over the canonical candidates and $g_{\rm curve}$ for the product
over the ones with two rational lifts. Given a monic degree-$k$
probe $Q$, one can compute
\[
   g_{\rm raw}\bmod Q,\qquad \gcd(g_{\rm curve},Q)
\]
in $\wtO(k+\sqrt{\tau k})$ field operations, with an additional
$O(\log p)$ factor allowed for modular powering. This applies
also on any dyadic subwindow, after updating the affine offset.

Choose a dyadic baby block length $L$ within a constant factor of
$\min(\tau,\sqrt{\tau k})$. Sort the $L$ residues
$b_i=(Ai+C)\bmod M$ and make their binary subproduct tree. For a
giant offset $d=ALj\bmod M$, the allowed outputs are $b_i+s$
with $s=d$ or $s=d-M$ and $0\le b_i+s<p$. These give two disjoint
intervals in the sorted list, each covered by $O(\log(2L))$ tree
nodes. For a selected node polynomial $P(Y)$ and shift $s$ we need
$P(X-s)\bmod Q(X)$. To batch all shifts for that node, compute
remainders of $P(Y)$ modulo $Q(Y+s)$ by a remainder tree and shift
each remainder back. The moduli need not be coprime. Multiply the
results modulo $Q$. Fast Taylor shifts follow from recursively
splitting polynomials at powers of two and precomputing the
corresponding powers of $X+s$.

The sum of the degrees of the baby polynomials over their tree is
$O(L\log(2L))$. There are $O((\tau/L)\log(2L))$ attached shifts.
This gives the bound $\wtO(L+k\tau/L+k)$. The baby/giant product
principle is classical. The two canonical intervals and batching
translated probe moduli handle the particular wrapping hashes here.
See \href{https://velusqrt.isogeny.org/velusqrt-20200616.pdf}
{Bernstein, De Feo, Leroux and Smith, Section 2}.
Finally put $D=\gcd(Q,g_{\rm raw}\bmod Q)$. It is squarefree
because the canonical abscissas are distinct. The second answer is
\[
  \gcd\big(D,(X^3+aX+c)^{(p-1)/2}-1\bmod D\big)
\]
with result one for $D=1$. This filters exactly the two-lift roots.
The gcd records common roots. The values of the omitted nonzero
factors would also be required to recover $g_{\rm curve}\bmod Q$.

\paragraph{What the discarded factors do to the target.}
On the odd prime-order subclass, there are no rational two-torsion
abscissas. Raw roots with no rational lift have their two lifts on
$E(\F_{p^2})$ exchanged by Frobenius, and those lifts satisfy
$\mathrm{Fr}(P)=-P$. At a rational signed target, Frobenius pairs
every nonrational signed solution without fixed points.
Consequently the raw and genuine target counts have identical
parity on every seed and on every restriction.

In fact at the tuned widths the probability of any extra solution
is $O(2^m\lambda/r)$ at each fixed eligible target.
Let $r_+=r$ and $r_-=2p+2-r$ be the sizes of the rational and
twist groups. For a nonempty set $J$ of twist slots the sum of its
points must lie in the intersection of the two groups, of size
$|E(\F_p)[2]|$. The complement then has a prescribed rational sum.
The all-twist case cannot reach $2T\ne O$. A fixed proper $J$
therefore contributes expected count at most
$|E(\F_p)[2]|r_-^{|J|-1}r_+^{m-|J|-1}\prod_i\rho_i$, where
$\rho_i=\tau_i/M$. Sum over $J$, use Hasse and
$\lambda=r^{m-1}\prod_i\rho_i$. Removal of candidates cannot
introduce a new tuple. On a clean isolated trial, the raw norm's
gradient differs by the nonzero product of the extra factors, so
coordinate ratios agree. The absolute scalar and first derivatives
are supplied by the three-block method with genuine inputs.

\section{Testing local recovery of the implicit resultant}
Write $A,B$ for the original projective half forms of degrees $f,h$.
All their roots are rational abscissas or infinity. Choose a fixed
nonlift $w$ or an element of $\F_{p^2}\setminus\F_p$ so neither
form vanishes there. The common change of coordinates
$(X_0,X_1)=(wY_0-Y_1,Y_0)$ has determinant one. Both transformed
affine polynomials have their full formal degrees and the leading
coefficient of the first is $a_f=A(w)$.

At $y=\infty$ write their quotient as a polynomial part plus
$\sum_{i\ge0}c_i y^{-i-1}$. Equivalently, with $z=1/y$, expand
$z^{f-h}B(w-z)/A(w-z)$. Let $H_f=(c_{i+j})_{0\le i,j<f}$. Then
\[
 \Res(A,B)=(-1)^{f(f-1)/2}a_f^{f+h}\det H_f,\qquad
 \operatorname{rank} H_f=f-\deg\gcd_{\rm proj}(A,B).
\]
For distinct finite roots $\alpha$ of the transformed first form,
the moments sum $\alpha^i B(\alpha)/A'(\alpha)$ in that chart.
The determinant of the resulting two Vandermonde matrices gives
the formula, which extends to repeated roots by polynomial identity
after clearing the unit leading coefficient. For rank, reduce the
proper part of the fraction to $P/D$ with $D$ monic. The form which
sends $(U,V)$ to the coefficient of $y^{-1}$ in $UVP/D$ is
nondegenerate on the quotient by $D$. Indeed, for a nonzero class $U$,
multiply the nonzero remainder of $UP$ by the power of $y$ that makes its degree $\deg D-1$.
The first $f$ monomials span this quotient. Its dimension is the
rank asserted above and the minimal order of a linear recurrence
for the negative-power moments.

On the actual independently hashed slots, the common degree is
at most the target count. Indeed the Sylvester matrix at the base
target has that corank, so its determinant after target deformation
vanishes to order at least that corank. This last fact follows by
constant row reduction and taking one local-parameter factor out
of each zero row. The determinant's order is exactly the ordinary
target count by the preceding norm identity. At an isolated target
the common degree is therefore exactly one. The series has exact
recurrence order $f-1$ there. At a fixed eligible target the mean
count is bounded, so at balanced widths the loss of rank is bounded
in probability. Local rational reconstruction from a short linear
recurrence does not bypass the half degree on these inputs.
This statement concerns the recurrence, not the minimum size of
arbitrary circuits computing its determinant.

A separate estimate concerns a test which picks candidate join
abscissas using the left slots and then checks for a right root.
Conditional on the left seeds, each specified signed point is the
required right sum with probability at most
$\prod_{i\in R}(r\rho_i)/r$ by the first moment bound. All but one
right coordinate can be chosen freely, and every eligible tuple survives
with probability $\prod_{i\in R}\rho_i$. A prescribed finite
abscissa allows two signs. Therefore $s$ candidates hit with
probability at most $2s\prod_{i\in R}(r\rho_i)/r$. The same holds
for a test that receives only failure until the first hit, or uses
moduli of total degree $s$ just to locate a common root in them.
A balanced split requires $s\ge 2^{b/2-o(b)}$ for constant success
in this interface. Arbitrary residue values can carry more
information than the hit predicate.

\section{A limit for arbitrary residue-value probes in a comparison model}
The following comparison grants the algorithm the entire second
polynomial for free and allows arbitrary computations on all
residues of the first. It clarifies which additional structure a
general local-probe reduction would need. Its inputs are fully
split random polynomials, distinct from the actual hash-seed law.

Put $d=\lfloor\sqrt p\rfloor$. Fix $Y\subset\F_p$ of size $d$
and $B=\prod_{y\in Y}(X-y)$. In the first distribution, $A$ has
$d$ independent uniform roots outside $Y$ (with replacement).
In the second, replace one such root by a uniform $y\in Y$.
Both forms are monic. The first resultant is nonzero and the
second has one simple common root. Even adaptive queries for
$A\bmod Q$ of total degree at most
$L=\lfloor\sqrt p/\log p\rfloor$ distinguish these distributions
with advantage $o(1)$. Repeated factors of $Q$ and probes in
extensions via irreducible moduli are allowed.

For a fixed transcript let $Q$ now denote the least common multiple
of all queried moduli and $\ell=\deg Q\le L$. Compare to a
reference oracle supplying consistent uniform unit residues. Each
such transcript has probability $1/|U_Q|$, where
$U_Q=(\F_p[X]/Q)^\times$, by successive uniform lifts. Conditional
on the free roots avoiding the roots of $Q$, a single root $a$ is
uniform on at least $p-d-\ell$ values. For a nontrivial character
$\chi$ of $U_Q$, extended by zero off the units,
\[
 \left|\E\chi(X-a\bmod Q)\right|
   \le\frac{\ell\sqrt p+d}{p-d-\ell}\le\frac2{\log p}
\]
for large $p$. Indeed the complete sum over $a$ is the linear
coefficient of the polynomial Dirichlet $L$-function for $\chi$.
It has degree at most $\ell-1$ and inverse roots of magnitude at
most $\sqrt p$ by the function-field Riemann hypothesis. The extra
Euler factors for an imprimitive character have inverse roots of
modulus one. Removing $Y$ costs at most $d$.
See \href{https://www.math.tau.ac.il/~rudnick/papers/KeatingRudnickIMRNfinal.pdf}
{Keating and Rudnick on Dirichlet $L$-functions}.

Multiplicative character orthogonality now gives pointwise relative
error at most $p^L(2/\log p)^{d-1}=o(1)$ from a uniform unit
residue after multiplying at least $d-1$ free factors. Multiplying
by a fixed unit from a planted root preserves this estimate.
The free roots avoid $Q$ with probability at least
$1-dL/(p-d)$, and the planted root avoids $Q$ with probability at
least $1-L/d$. Thus in either distribution each unit transcript
has probability at least $(1-o(1))/|U_Q|$, uniformly over the
choice of transcript. Summing the common mass with the reference
distribution proves the claim, including adaptive probes and
randomized algorithms. The argument only uses residues of a hidden
form. An algorithm that exploits the actual public seeds directly
is outside this comparison theorem.

\section{End-to-end result and finite evidence}
The project's conditional end-to-end bound is still
\[
  \log_2\E T_{\rm solve}\le (26.2933775156+o(1))
                 \frac{b}{(\log_2 b)^{1/4}}
\]
under its inherited contraction hypotheses and finite-packet
assumption. Making both half polynomials implicit reduced working
space. The section evaluations still cost an exponential in $b$,
and the faster genuine-slot gcd probe does not compute the
normalized resultant of two aggregated halves. A faster simultaneous
evaluation using the internal slot descriptions remains open.

Run \texttt{python -B research/implicit\_resultant/verify.py}.
The comparison uses multiplication-matrix cofactors for 567
full-gradient checks over all 63 dyadic boxes. It also checks
multiple zeros, two-torsion, and 45 independent Hankel determinants
and ranks. On the order-65419 fixture the old degrees were 56 and
400. The three-block evaluator used degrees 80, 14 and 10, and no
temporary chord section exceeded degree 28. It reproduced the zero
norm, target derivative 43446, and additive derivatives
16953, 13292, 47523, 12119 and 9163 in $\F_{65521}$. The explicit
half forms are used only in separate verification comparisons.
The degree-56 denominator had moment rank 55 at this isolated root.
There were 2052 raw-remainder and genuine-gcd comparisons up to
width 16384, including repeated-root probes of degree 16. All 27 raw/curve parity comparisons agreed, with six
strict differences in integer multiplicity. The time bounds above come from counting operations.
\end{document}
